% Einheit 4 – Graph und Ableitungsgraph
\setcounter{aufgabe}{30}
\hypertarget{e4}{}\einheitenkopf{Einheit 4 von 5 · Graph und Ableitungsgraph}
\zweigzeile{Hier lernst du, aus dem Graphen einer Funktion den Graphen ihrer Ableitung zu entwickeln und umgekehrt · kennst du seit Q1 (Klasse 11), je nach Kursplan bis Q2 · Abitur GK · baut auf: Monotonie (Einheit 1), Extrempunkte (Einheit 2), Krümmung und Wendepunkte (Einheit 3), Ableitung als Steigung deuten, Graphen grob skizzieren}

\begin{aufgabe}{Ich kann am Graphen von $f$ ablesen, wo $f'$ positiv, negativ oder null ist. Die Abbildungen zeigen die Graphen von $f_1$ und $f_2$; in den Graphen von $f_2$ ist eine Tangente eingezeichnet.}
Beispiel: Der Graph steigt bis zum Hochpunkt bei $x = -1$, fällt bis zum Tiefpunkt bei $x = 2$ und steigt dann wieder. \\
Also: $f'(x) = 0$ bei $x = -1$ und $x = 2$; \ $f'(x) > 0$ für $x < -1$ und $x > 2$; \ $f'(x) < 0$ für $-1 < x < 2$.

\begin{ksys}[xmin=-3,xmax=4,ymin=-3,ymax=4,ablesen]
\funktionab{0.25*\x^4-2*\x^2+1}{f_1}{-3}{3}
\end{ksys}\ksysabstand\begin{ksys}[xmin=-2,xmax=4,ymin=-3,ymax=5,ablesen]
\funktionab{0.5*\x^3-1.5*\x^2+2}{f_2}{-1.5}{3.5}
\tangentean*{0.5*\x^3-1.5*\x^2+2}{1}{t}
\end{ksys}

\begin{teile}
\teil Gib die Stellen an, an denen $f_1'(x) = 0$ ist. \leerfeld
\teil Vorzeichen von $f_1'$ in den Abschnitten, von links nach rechts: \leerfeld
\teil Gib die Stellen an, an denen $f_2'(x) = 0$ ist. \leerfeld
\teil Gib das Vorzeichen von $f_2'(-1)$, $f_2'(1)$ und $f_2'(3)$ an. \leerfeld
\teil Markiere den Wendepunkt von $f_2$, begründe, dass die Steigung dort negativ ist, und lies $f_2'$ an der Wendestelle an der Tangente ab.
\teil Markiere auf dem Graphen von $f_2$ einen Punkt, in dem $f_2(x) > 0$, $f_2'(x) < 0$ und $f_2''(x) > 0$ gilt.
\end{teile}
\end{aufgabe}

\begin{aufgabe}{Ich kann den Graphen der Ableitung zum Graphen von $f$ zeichnen. Oben ist der Graph von $f(x) = \frac{1}{3}x^3 - x^2$ gezeichnet.}
\begin{teile}
\teil Bilde $f'(x)$ und fülle die Wertetabelle aus.
\end{teile}
\wertetabelle{x}{f'(x)}{-1,0,1,2,3}

\ableitungspaar[xmin=-1,xmax=3,ymin=-2,ymax=1,ymin2=-2,ymax2=4]{\x^3/3-\x^2}{f}{}{f'}

\begin{teile}
\teil Zeichne den Graphen von $f'$ in das untere Koordinatensystem.
\teil Begründe am Bild, warum die Nullstellen von $f'$ genau unter dem Hoch- und dem Tiefpunkt von $f$ liegen.
\end{teile}
\end{aufgabe}

\begin{aufgabe}{Ich kann einen möglichen Graphen zu vorgegebenen Eigenschaften skizzieren. Eine ganzrationale Funktion $f$ dritten Grades hat den Hochpunkt $H(-2 \mid 2)$, den Tiefpunkt $T(2 \mid -3)$ und die Wendestelle $x = 0$.}
\begin{teile}
\teil Skizziere einen möglichen Graphen von $f$ in das obere Koordinatensystem.
\teil Skizziere dazu den Graphen von $f'$ in das untere; beachte, dass die Wendestelle von $f$ eine Extremstelle von $f'$ ist.
\end{teile}

\ableitungspaar[xmin=-4,xmax=4,ymin=-4,ymax=3,ymin2=-3,ymax2=3]{}{f}{}{f'}

\begin{teile}
\teil Begründe, welchen Grad eine ganzrationale Funktion mindestens hat, deren Graph zwei Tiefpunkte und einen Hochpunkt hat.
\end{teile}
\end{aufgabe}

\begin{aufgabe}{Ich kann aus dem Graphen einer Differenzfunktion die Lage zweier Graphen ablesen. Die Abbildung zeigt den Graphen der Funktion $d$ mit $d(x) = f(x) - g(x)$.}

\begin{ksys}[xmin=-2,xmax=4,ymin=-3,ymax=3,ablesen]
\funktionab{-0.5*\x^2+\x+1.5}{d}{-2}{4}
\end{ksys}

\begin{teile}
\teil Gib an, für welche $x$ der Graph von $f$ oberhalb des Graphen von $g$ liegt. \leerfeld
\teil Gib die Stellen an, an denen sich die Graphen von $f$ und $g$ schneiden. \leerfeld
\teil Gib an, an welcher Stelle zwischen $x = -1$ und $x = 3$ der senkrechte Abstand der beiden Graphen am größten ist und wie groß er dort ist.
\end{teile}
\end{aufgabe}

\begin{aufgabe}{Ich kann Aussagen über einen Graphen und seine Ableitung beurteilen. Die Abbildung zeigt den Graphen von $f(x) = -0{,}5x^3 + 1{,}5x^2 + 1$; es ist $f'(x) = -1{,}5x^2 + 3x$.}

\begin{ksys}[xmin=-2,xmax=4,ymin=-3,ymax=5,ablesen]
\funktionab{-0.5*\x^3+1.5*\x^2+1}{f}{-1.3}{3.5}
\end{ksys}

\begin{teile}
\teil Beurteile die beiden Aussagen: (1) „Die Normale im Wendepunkt des Graphen von $f$ ist fallend und hat die Steigung $-\frac{2}{3}$.“ (2) „Die Ableitung $f'$ nimmt jeden Wert zwischen $-3$ und $3$ an.“ (Abitur 2024)
\teil Beurteile die Aussage „Jede Tangente an den Graphen von $f$ hat mit dem Graphen genau zwei gemeinsame Punkte.“ Zeichne dazu passende Tangenten in die Abbildung ein. (Abitur 2023)
\end{teile}
\end{aufgabe}

\begin{aufgabe}{Ich finde den Fehler beim Lesen eines Ableitungsgraphen. Die Abbildung zeigt den Graphen der Ableitung $f'$ einer Funktion $f$. Ben schreibt: „Der Graph hat bei $x = 1$ seinen Tiefpunkt, also hat $f$ bei $x = 1$ einen Tiefpunkt.“}

\begin{ksys}[xmin=-2,xmax=4,ymin=-5,ymax=5,ablesen]
\funktionab{\x^2-2*\x-3}{f'}{-2}{4}
\end{ksys}

\begin{teile}
\teil Finde den Fehler, benenne ihn und korrigiere die Aussage.
\teil Bestimme am Graphen von $f'$ die Extremstellen von $f$ und ihre Art.
\end{teile}
\end{aufgabe}

\begin{aufgabe}{Ich kann Zusammenhänge zwischen $f$ und $f'$ begründen.}
\begin{teile}
\teil Begründe, warum der Graph von $f'$ bei einer ganzrationalen Funktion $f$ einen Grad niedriger ist als der von $f$.
\teil Begründe: Hat der Graph von $f$ an der Stelle $x_0$ einen Wendepunkt, so hat der Graph von $f'$ dort einen Hoch- oder Tiefpunkt.
\teil $f$ und $g$ sind ganzrationale Funktionen, deren Graphen auf dem Intervall $[0;\,4]$ oberhalb der $x$-Achse verlaufen. Die Flächen, die die beiden Graphen dort mit der $x$-Achse einschließen, sind gleich groß. Begründe, dass sich die Graphen von $f$ und $g$ im Intervall $[0;\,4]$ mindestens einmal schneiden.
\end{teile}
\end{aufgabe}

\begin{aufgabe}{Ich kann aus dem Graphen einer Änderungsrate den Verlauf eines Bestands ablesen. Die Abbildung zeigt, wie schnell sich der Wasserstand eines Stausees ändert: $w'(t)$ in Zentimetern pro Stunde, $t$ in Stunden nach Mitternacht.}

\begin{ksys}[xmin=0,xmax=12,ymin=-6,ymax=5,ablesen,xlabel=$t$ in h,ylabel=$w'(t)$ in cm/h]
\funktionab{-0.25*\x^2+3*\x-5}{w'}{0}{12}
\end{ksys}

\begin{teile}
\teil Gib an, in welchem Zeitraum der Wasserstand steigt. \leerfeld
\teil Gib an, wann der Wasserstand einen Tiefpunkt und wann er einen Hochpunkt erreicht.
\teil Gib an, wann der Wasserstand am schnellsten steigt und wie schnell.
\end{teile}
\end{aufgabe}
